% ============================================================================
% ARBEITSBLATT DS 9 - Gruppe B (B1): Repaso y Preparación
% ============================================================================
% Spanisch Spätbeginner - Einstimmung vor Beobachtungsstunde
% Fokus: Wiederholung ser/estar, tener, hay que, Personenbeschreibung (erweitertes Niveau)
% ============================================================================

\documentclass[11pt,a4paper]{article}

\usepackage[utf8]{inputenc}
\usepackage[T1]{fontenc}
\usepackage[ngerman]{babel}
\usepackage[left=2cm,right=2cm,top=1.8cm,bottom=1.5cm]{geometry}
\usepackage{helvet}
\usepackage[table]{xcolor}
\usepackage{tabularx}
\usepackage{array}
\usepackage{enumitem}
\usepackage{tcolorbox}
\usepackage{fancyhdr}
\usepackage{lastpage}
\usepackage{amssymb}

\renewcommand{\familydefault}{\sfdefault}

\definecolor{spanisch}{HTML}{c41e3a}
\definecolor{gruppeB}{HTML}{1565c0}
\definecolor{grau}{HTML}{666666}
\definecolor{hellgrau}{HTML}{f5f5f5}

\pagestyle{fancy}
\fancyhf{}
\fancyhead[L]{\small\textcolor{grau}{Spanisch Spätbeginner -- Gruppe B (B1)}}
\fancyhead[R]{\small\textcolor{gruppeB}{AB-09-B}}
\fancyfoot[C]{\small\thepage/\pageref{LastPage}}
\renewcommand{\headrulewidth}{0.5pt}

\newcommand{\luecke}[1]{\underline{\hspace{#1}}}
\newcommand{\punktebox}[1]{\hfill\fbox{\small \_\_/#1}}
\newcommand{\es}[1]{\textit{#1}}

\newtcolorbox{merkkasten}[1][]{
    colback=white,
    colframe=gruppeB,
    fonttitle=\bfseries\color{gruppeB},
    title=#1,
    boxrule=1.5pt,
    arc=0pt,
    left=5pt, right=5pt, top=5pt, bottom=5pt
}

\newtcolorbox{hilfekasten}{
    colback=hellgrau,
    colframe=grau,
    boxrule=0.5pt,
    arc=0pt,
    left=5pt, right=5pt, top=3pt, bottom=3pt
}

\newcounter{aufgabe}
\newcommand{\aufgabe}[2]{%
    \stepcounter{aufgabe}%
    \vspace{0.8em}%
    \noindent\textbf{\theaufgabe.} #1 \punktebox{#2}%
    \vspace{0.3em}%
}

\begin{document}

% ============================================================================
% KOPFBEREICH
% ============================================================================
\noindent
\begin{tabularx}{\textwidth}{@{}Xr@{}}
    {\Large\textbf{\textcolor{gruppeB}{Repaso avanzado: Describir personas}}} & Name: \luecke{5cm} \\[0.3em]
    {\small DS 9 -- Wiederholung zur Vorbereitung} & Datum: \luecke{3cm} \\
\end{tabularx}

\vspace{0.3em}
\hrule
\vspace{0.8em}

% ============================================================================
% MERKKASTEN 1: ser vs. estar - erweitert
% ============================================================================
\begin{merkkasten}[Kasten 1: \es{ser} vs. \es{estar} -- Unterschiede und Nuancen]
\vspace{0.2em}

\begin{tabularx}{\textwidth}{@{}|p{2cm}|X|X|@{}}
\hline
\rowcolor{hellgrau}
\textbf{Verb} & \textbf{Verwendung} & \textbf{Beispiele} \\
\hline
\textbf{ser} & \luecke{3.5cm} & \es{Soy inteligente.} / \es{Es de Madrid.} \\[0.5em]
\hline
\textbf{estar} & \luecke{3.5cm} & \es{Estoy cansado.} / \es{Está en casa.} \\[0.5em]
\hline
\end{tabularx}

\vspace{0.5em}
\textbf{Achtung!} Manche Adjektive ändern ihre Bedeutung:

\vspace{0.3em}
\begin{tabular}{ll}
\es{ser listo} = \luecke{3cm} & \es{estar listo} = \luecke{3cm} \\
\es{ser aburrido} = \luecke{3cm} & \es{estar aburrido} = \luecke{3cm} \\
\end{tabular}
\vspace{0.3em}
\end{merkkasten}

\vspace{0.8em}

% ============================================================================
% MERKKASTEN 2: Konnektoren
% ============================================================================
\begin{merkkasten}[Kasten 2: Konnektoren für komplexe Beschreibungen]
\vspace{0.2em}

\begin{tabularx}{\textwidth}{@{}X|X|X@{}}
\es{también} = auch & \es{pero} = aber & \es{porque} = weil \\
\es{además} = außerdem & \es{sin embargo} = jedoch & \es{ya que} = da, weil \\
\end{tabularx}

\vspace{0.3em}
\textbf{Beispiel:} \es{Mi amiga es inteligente y \luecke{2cm} muy simpática, \luecke{2cm} a veces está cansada \luecke{2cm} trabaja mucho.}
\vspace{0.3em}
\end{merkkasten}

\vspace{1em}

% ============================================================================
% AUFGABEN
% ============================================================================

\aufgabe{Ergänze \es{ser} oder \es{estar} in der richtigen Form. Achte auf den Bedeutungsunterschied!}{6}

\begin{enumerate}[label=\alph*), leftmargin=2em]
    \item Mi profesor \luecke{2cm} muy aburrido. (= Er ist langweilig.)
    \item Hoy \luecke{2cm} aburrida porque no hay nada que hacer. (= Ich langweile mich.)
    \item Los estudiantes \luecke{2cm} listos para el examen. (= Sie sind bereit.)
    \item Tu hermana \luecke{2cm} muy lista. (= Sie ist sehr klug.)
    \item Normalmente \luecke{2cm} tranquilo, pero hoy \luecke{2cm} nervioso.
    \item ¿Dónde \luecke{2cm} María? -- \luecke{2cm} en la biblioteca.
\end{enumerate}

\vspace{0.5em}

\aufgabe{Wortschatz erweitert: Adjektive und Antonyme.}{6}

\begin{tabularx}{\textwidth}{@{}|X|X|X|X|@{}}
\hline
\rowcolor{hellgrau}
\textbf{Adjektiv} & \textbf{Übersetzung} & \textbf{Antonym} & \textbf{Übersetzung} \\
\hline
simpático/-a & \luecke{2cm} & antipático/-a & \luecke{2cm} \\[0.5em]
\hline
trabajador/-a & \luecke{2cm} & perezoso/-a & \luecke{2cm} \\[0.5em]
\hline
optimista & \luecke{2cm} & pesimista & \luecke{2cm} \\[0.5em]
\hline
\end{tabularx}

\vspace{1em}

\aufgabe{Verbinde die Sätze mit passenden Konnektoren.}{6}

\begin{hilfekasten}
\textbf{Konnektoren:} \es{también} | \es{pero} | \es{porque} | \es{además} | \es{sin embargo} | \es{ya que}
\end{hilfekasten}

\vspace{0.3em}
\begin{enumerate}[label=\alph*), leftmargin=2em]
    \item Mi amigo es simpático. Es inteligente.

    $\rightarrow$ \luecke{12cm}

    \item Quiero ser médico. Me gusta ayudar a la gente.

    $\rightarrow$ \luecke{12cm}

    \item Es muy trabajadora. A veces está cansada.

    $\rightarrow$ \luecke{12cm}
\end{enumerate}

\vspace{0.5em}

\aufgabe{Berufe und Eigenschaften: Schreibe vollständige Sätze mit \es{hay que} und Begründung.}{8}

\vspace{0.3cm}
\textbf{Beispiel:} \textit{médico/-a} $\rightarrow$ \es{Para ser médico hay que ser paciente y responsable, ya que los pacientes necesitan mucha atención.}

\vspace{0.4cm}
a) \textbf{profesor/-a:}

\luecke{14cm}

\vspace{0.4cm}
b) \textbf{artista:}

\luecke{14cm}

\vspace{0.4cm}
c) \textbf{bombero/-a:}

\luecke{14cm}

\vspace{0.4cm}
d) \textbf{veterinario/-a:}

\luecke{14cm}

\vspace{1em}

\aufgabe{Textproduktion: Beschreibe eine Person aus deiner Klasse (6--8 Sätze).}{9}

\begin{hilfekasten}
\textbf{Verwende:} \es{ser/estar}, mindestens 3 verschiedene Konnektoren (\es{también, pero, porque, además}), \es{tener + años}, Adjektive mit Steigerung (\es{muy, bastante, un poco})
\end{hilfekasten}

\vspace{0.5em}
\textbf{Mi compañero/-a de clase}

\vspace{0.3em}
\begin{tabularx}{\textwidth}{@{}|X|@{}}
\hline
\\[0.5em]
\luecke{14cm} \\[0.8em]
\luecke{14cm} \\[0.8em]
\luecke{14cm} \\[0.8em]
\luecke{14cm} \\[0.8em]
\luecke{14cm} \\[0.8em]
\luecke{14cm} \\[0.8em]
\luecke{14cm} \\[0.8em]
\hline
\end{tabularx}

\vspace{1em}
\hrule
\vspace{0.3em}
\begin{flushright}
\textbf{Gesamt: \luecke{1cm} / 35 Punkte}
\end{flushright}

\end{document}
